\documentclass[11pt]{amsart}
\usepackage{calc,amssymb,amsmath,ulem,stmaryrd,fullpage}
\usepackage{enumerate}
\usepackage{alltt}
\RequirePackage[dvipsnames,usenames]{color}
\let\mffont=\sf
\normalem
%\input{kmacros3.sty}
\input{xy}
\xyoption{all}
\usepackage{hyperref}
\usepackage{graphicx}
\usepackage[all,cmtip]{xy}
\usepackage{verbatim}
\usepackage[left=0.95in,top=0.95in,right=0.95in,bottom=0.95in]{geometry}
\newcommand{\tauCohomology}{T}
\newcommand{\FCohomology}{S}


\theoremstyle{theorem}
%\newtheorem*{mainthm*}{Main Theorem}

\newcommand{\note}[1]{\marginpar{\sffamily\tiny #1}}

\def\todo#1{\textcolor{Mahogany}%
{\sffamily\footnotesize\newline{\color{Mahogany}\fbox{\parbox{\textwidth-15pt}{#1}}}\newline}}

\renewcommand{\O}{\mathcal O}
\newcommand{\ie}{{\itshape i.e.} }
\newcommand{\roundup}[1]{\lceil #1 \rceil}
\newcommand{\rounddown}[1]{\lfloor #1 \rfloor}
\renewcommand{\to}[1][]{\xrightarrow{\ #1\ }}
\newcommand{\onto}[1][]{\protect{\xrightarrow{\ #1\ }\hspace{-0.8em}\rightarrow}}
\newcommand{\into}[1][]{\lhook \joinrel \xrightarrow{\ #1\ }}
%\newcommand{DBDual}[1]{{\underline{\omega}_{#1}}}
\newcommand{\DBDual}[1]{{{\uuline \omega}{}^{\mydot}_{#1}}}
\newcommand{\Tt}{{\mathfrak{T}}}
%\newcommand{\bn}{{\mathfrak{n}} }
\newcommand{\sol}[1]{ \vskip 6pt{\color{blue} {\bf Solution:} #1} \vskip 6pt}

\newcommand{\bR}{{\mathbb{R}}}
\newcommand{\va}{{\vec{a}}}
\newcommand{\vb}{{\vec{b}}}
\newcommand{\vzero}{{\vec{0}}}
\newcommand{\vi}{{\vec{i}}}
\newcommand{\vj}{{\vec{j}}}
\newcommand{\vk}{{\vec{k}}}
\newcommand{\vh}{{\vec{h}}}
\newcommand{\vr}{{\vec{r}}}
\newcommand{\vu}{{\vec{u}}}
\newcommand{\vv}{{\vec{v}}}
\newcommand{\vw}{{\vec{w}}}
\newcommand{\vx}{{\vec{x}}}
\newcommand{\vy}{{\vec{y}}}
\newcommand{\vz}{{\vec{z}}}
\newcommand{\vT}{{\vec{T}}}
\newcommand{\vN}{{\vec{N}}}
\newcommand{\vF}{{\vec{F}}}
\newcommand{\vG}{{\vec{G}}}
\newcommand{\bF}{\mathbb{F}}
\newcommand{\bQ}{\mathbb{Q}}
\newcommand{\bZ}{\mathbb{Z}}
\newcommand{\bC}{\mathbb{C}}
\newcommand{\Gal}{\textnormal{Gal}}
\newcommand{\Aut}{\textnormal{Aut}}
\newcommand{\Emb}{\textnormal{Emb}}
\newcommand{\Tr}{\textnormal{Tr}}
\newcommand{\N}{\textnormal{N}}
\newcommand{\Hom}{\textnormal{Hom}}
\newcommand{\Ext}{\textnormal{Ext}}
\newcommand{\Spec}{\textnormal{Spec}}
\newcommand{\fram}{\mathfrak{m}}
\newcommand{\id}{\textnormal{id}}
\newcommand{\init}{\textnormal{in}}
\newcommand{\lex}{\textnormal{lex}}
\newcommand{\grevlex}{\textnormal{grevlex}}

\begin{document}
\title{HW \#8 -- Math 6320\\Spring 2015}
\author{Due: Tuesday April 28th}
\maketitle

\begin{enumerate}
\item Suppose $S = k[x_1, \ldots, x_r]$ and $f \in S$.  If $\init_{\lex}(f) \in k[x_s, \ldots, x_r]$ for some $s$, then show that $f \in k[x_s, \ldots, x_r]$.
\item Suppose $S = k[x_1, \ldots, x_r]$ and $f \in S$.  If $f$ is homogeneous with $\init_{\grevlex}(f) \in \langle x_s, \ldots, x_r \rangle$, for some $s$, then show that $f \in \langle x_s, \ldots, x_r \rangle$.
\item Say $S = k[x_1, \ldots, x_r]$ and that $F$ is a finitely generated free $S$-module with basis and a fixed monomial order $>$.  Suppose that $g_1, \ldots, g_t \in M \subseteq F$.  Suppose that $\init(g_1), \ldots, \init(g_t)$ generated $\init(M)$, then show that $g_1, \ldots, g_t$ generate $M$ and hence is a Gr\"obner basis for $M$.  
\item Suppose that $R = \bR[x]$, $M$ and $N$ are finitely generated $R$-modules and $L = R/\langle x^2 + 1 \rangle \oplus R^2$.  Suppose that $L \oplus M \cong L \oplus N$.  Prove that $M \cong N$.
\item Find an ideal $I$ in the ring $A = \bZ[x]$ such that $A/I$ has exactly three prime ideals. Identify the
ideals and justify your assertion.
\item Let $G$ be a finite group of order $p^n$, where $p$ is prime and $n\geq 1$. Suppose $G$
acts on a finite set $S$. Let $S'$ be the subset of $S$ consisting of elements fixed by $G$:
$$
S'=\{x\in S\ |\ gx=x\text{ for all }g\in G\}.
$$
Prove that the order of $S'$ is congruent to the order of $S$ modulo $p$.
\item Consider $f = x^6 + 3 \in \bQ[x]$  and let $\alpha \in \bC$ be a root of $f$.  Set $E = \bQ[\alpha]$.  
\begin{itemize}
\item[(a)]  Show that $E$ contains a primitive $6$th root of unity.  
\item[(b)]  Show that $E$ is Galois over $\bQ$.  
\item[(c)]  How many intermediate fields $\bQ \subseteq F \subseteq E$ are there with $|F : \bQ| = 3$
\end{itemize}
\end{enumerate}
\end{document}

